\documentclass[10pt,a4paper]{article}
\usepackage[utf8]{inputenc}
\usepackage[T5]{fontenc}
\usepackage[vietnamese]{babel}
\usepackage[margin=1.4cm,top=0.85cm,bottom=1.05cm,includehead,includefoot,headheight=20pt,headsep=6pt,footskip=16pt]{geometry}
\usepackage{amsmath,amssymb,array,tabularx,enumitem,fancyhdr,tikz,xcolor,needspace,multicol,colortbl,microtype}
\usepackage{tgheros}
\usetikzlibrary{calc,angles,quotes,arrows.meta,babel}

% =========================================================================
% MINIMALIST SLATE-NAVY DESIGN SYSTEM (UI-STYLING / PRINT EDITORIAL)
% Nguyên tắc: Sạch, thoáng, tương phản cao, 1 màu nhấn Navy duy nhất,
% viền mảnh, không tô nền sặc sỡ -> dễ đọc cả trên màn hình lẫn bản in đen trắng.
% =========================================================================
\definecolor{ink}{HTML}{0F172A}       % Slate 900 - Chữ tiêu đề / nội dung chính
\definecolor{navy}{HTML}{1E3A5F}      % Academic Navy - Màu nhấn duy nhất
\definecolor{slateBody}{HTML}{334155} % Slate 700 - Nhãn phụ
\definecolor{muted}{HTML}{64748B}     % Slate 500 - Chú thích, quote, số trang
\definecolor{hairline}{HTML}{CBD5E1}  % Slate 300 - Đường kẻ phân tách mảnh
\definecolor{surface}{HTML}{F8FAFC}   % Slate 50  - Nền xám rất nhẹ (ô đáp số, header bảng)

% ---- Bảng màu riêng cho hình vẽ TikZ (tách biệt rõ nét các đối tượng/vectơ) ----
\definecolor{figEdge}{HTML}{1E3A5F}      % Viền hình chính (Navy đậm)
\definecolor{figFill}{HTML}{EFF6FF}      % Tô nền hình nhẹ (Sky 50)
\definecolor{figPrimary}{HTML}{0284C7}   % Vectơ / đường phụ 1 (Xanh dương sáng)
\definecolor{figSecondary}{HTML}{D97706} % Vectơ / đường phụ 2 (Cam hổ phách)
\definecolor{figHighlight}{HTML}{DC2626} % Góc vuông, góc xen giữa, điểm nhấn (Đỏ san hô)
\definecolor{figAxis}{HTML}{64748B}      % Trục tọa độ, đường gióng (Xám slate)

% ---- Thông số đề (chỉnh nhanh tại đây) ----
\newcommand{\examSeries}{LÀM TOÁN ĐỂ HẠNH PHÚC}
\newcommand{\examKind}{ĐỀ ÔN TẬP SỐ 01}
\newcommand{\examTitle}{ĐỀ ÔN TẬP TOÁN 10 — CHƯƠNG IV. VECTƠ}
\newcommand{\examYear}{NĂM HỌC 2026 -- 2027}
\newcommand{\examMeta}{Môn: TOÁN -- Khối lớp 10}
\newcommand{\examTime}{Thời gian làm bài: 90 phút (không kể thời gian phát đề)}
\newcommand{\examCode}{101}
\newcommand{\examAuthor}{Dương Gia Bảo}
\newcommand{\examPages}{1}
\newcommand{\ansPages}{1}

\setlength{\parindent}{0pt}
\setlength{\parskip}{0pt}
\setlength{\columnseprule}{0.4pt}
\def\columnseprulecolor{\color{hairline}}

\newcommand{\headerQuote}{%
  \ifcase\value{page}\relax
    \or Math gives us hope that every problem has a solution%
    \or Math + You = Champion%
    \or Where Frustration Becomes Inspiration%
    \or Turning Confusion into Clarity%
    \or Don't be an X, be a Y of mathematics%
    \or Don't get mad. Learn to add%
    \else Math gives us hope that every problem has a solution%
  \fi
}

\pagestyle{fancy}
\fancyhf{}
\lhead{\footnotesize\textsf{\textit{\color{muted}\headerQuote}}}
\rhead{\footnotesize\textsf{\color{muted}Người biên soạn:\ \color{ink}\textbf{\examAuthor}}}
\lfoot{\footnotesize\textsf{\color{muted}Toán 10 KNTT}}
\rfoot{\footnotesize\textsf{\color{muted}Trang \color{ink}\textbf{\thepage/\examPages}}}
\renewcommand{\headrulewidth}{0.4pt}
\renewcommand{\headrule}{\hbox to\headwidth{\color{hairline}\leaders\hrule height \headrulewidth\hfill}}
\renewcommand{\footrulewidth}{0pt}

\fancypagestyle{dapan}{
  \fancyhf{}
  \lhead{\footnotesize\textsf{\textit{\color{muted}\headerQuote}}}
  \rhead{\footnotesize\textsf{\color{muted}Người biên soạn:\ \color{ink}\textbf{\examAuthor}}}
  \lfoot{\footnotesize\textsf{\textit{\color{muted}Đáp án \& Hướng dẫn giải — \examKind}}}
  \rfoot{\footnotesize\textsf{\color{muted}Trang \color{ink}\textbf{\thepage/\ansPages}}}
  \renewcommand{\headrulewidth}{0.4pt}
  \renewcommand{\headrule}{\hbox to\headwidth{\color{hairline}\leaders\hrule height \headrulewidth\hfill}}
  \renewcommand{\footrulewidth}{0pt}
}

% ---- Macro trình bày tối giản ----
\newcommand{\parttitle}[1]{%
  \par\vspace{6pt}\needspace{3\baselineskip}%
  \noindent{\color{navy}\rule[-2.5pt]{2pt}{11.5pt}}\hspace{5pt}%
  {\sffamily\normalsize\bfseries\color{ink}#1}%
  \par\vspace{1.5pt}{\color{hairline}\hrule height 0.45pt}\vspace{3pt}\nopagebreak%
}
\newcommand{\partnote}[1]{{\footnotesize\color{muted}\textit{#1}}\par\vspace{2pt}}

\newcommand{\four}[4]{\par\vspace{1.5pt}\noindent\begin{tabularx}{\linewidth}{@{}*{4}{>{\raggedright\arraybackslash}X}@{}}\textsf{\textbf{\color{ink}A.}} #1&\textsf{\textbf{\color{ink}B.}} #2&\textsf{\textbf{\color{ink}C.}} #3&\textsf{\textbf{\color{ink}D.}} #4\end{tabularx}}
\newcommand{\two}[4]{\par\vspace{1.5pt}\noindent\begin{tabularx}{\linewidth}{@{}*{2}{>{\raggedright\arraybackslash}X}@{}}\textsf{\textbf{\color{ink}A.}} #1&\textsf{\textbf{\color{ink}B.}} #2\\[4pt]\textsf{\textbf{\color{ink}C.}} #3&\textsf{\textbf{\color{ink}D.}} #4\end{tabularx}}
\newcommand{\one}[4]{\par\vspace{1.5pt}\noindent\begin{tabularx}{\linewidth}{@{}>{\raggedright\arraybackslash}X@{}}\textsf{\textbf{\color{ink}A.}} #1\\[4pt]\textsf{\textbf{\color{ink}B.}} #2\\[4pt]\textsf{\textbf{\color{ink}C.}} #3\\[4pt]\textsf{\textbf{\color{ink}D.}} #4\end{tabularx}}

\newcommand{\answerbox}{\unskip{\nobreak\hfill\penalty50\hskip1em\hbox{}\nobreak\hfill\mbox{\textsf{\footnotesize\color{muted}Đáp số:}\ \tikz[baseline=-3pt]\node[draw=hairline!70!ink,line width=0.55pt,fill=surface,rounded corners=2pt,minimum width=1.85cm,minimum height=13.5pt,inner sep=0pt]{};}}}
\newcommand{\essayitem}[2]{\par\vspace{4pt}\needspace{3\baselineskip}\textsf{\textbf{\color{navy}#1.}}\ \textsf{\footnotesize\color{muted}(#2)}\ }

\setlist[enumerate,1]{label=\textsf{\textbf{\color{navy}Câu \arabic*.}},leftmargin=1.18cm,labelsep=.15cm,itemsep=3pt,topsep=2pt,parsep=0pt}
\setlist[enumerate,2]{label=\textsf{\color{slateBody}\alph*)},leftmargin=.55cm,itemsep=2pt,topsep=1pt,parsep=0pt}
\newcolumntype{C}[1]{>{\centering\arraybackslash}m{#1}}

\begin{document}

% ==================== HEADER ĐỀ THI TỐI GIẢN ====================
{\color{navy}\hrule height 1.2pt}\vspace{6pt}
\noindent
\begin{minipage}[c]{0.31\textwidth}
    \centering
    {\footnotesize\sffamily\color{muted}\examSeries}\\[2pt]
    {\color{hairline}\rule{2.4cm}{0.5pt}}\\[3pt]
    {\normalsize\sffamily\bfseries\color{ink}\examKind}
\end{minipage}\hfill
\begin{minipage}[c]{0.52\textwidth}
    \centering
    {\normalsize\sffamily\bfseries\color{ink}\examTitle}\\[2pt]
    {\small\sffamily\color{slateBody}\examYear\ \textbullet\ \examMeta}\\[1pt]
    {\footnotesize\textit{\color{muted}\examTime}}
\end{minipage}\hfill
\begin{minipage}[c]{0.14\textwidth}
    \raggedleft
    \tikz\node[draw=navy,line width=0.65pt,fill=surface,rounded corners=2.5pt,inner xsep=8pt,inner ysep=4pt,align=center]{%
        {\scriptsize\sffamily\color{muted}MÃ ĐỀ}\\[-1pt]
        {\normalsize\sffamily\bfseries\color{navy}\examCode}%
    };
\end{minipage}

\vspace{6pt}{\color{hairline}\hrule height 0.45pt}\vspace{5pt}
\noindent\textsf{\small\color{ink}\textbf{Họ và tên thí sinh:}}\dotfill\quad\textsf{\small\color{ink}\textbf{SBD:}}\makebox[3.2cm]{\dotfill}\quad\textsf{\small\color{ink}\textbf{Lớp:}}\makebox[1.8cm]{\dotfill}
\vspace{4pt}{\color{hairline}\hrule height 0.45pt}\vspace{2pt}

% ==================== PHẦN I: TRẮC NGHIỆM 4 PHƯƠNG ÁN ====================
\parttitle{PHẦN I. Câu trắc nghiệm nhiều phương án lựa chọn}
\partnote{Thí sinh trả lời từ câu 1 đến câu 4. Mỗi câu hỏi thí sinh chỉ chọn một phương án.}
\begin{enumerate}
\item Cho hình bình hành $ABCD$ tâm $O$. Khẳng định nào sau đây \textbf{đúng}?
\four{$\vec{AB} = \vec{CD}$.}{$\vec{OA} = \vec{OC}$.}{$\vec{AB} = \vec{DC}$.}{$\vec{AD} = \vec{CB}$.}

\item Trong mặt phẳng tọa độ $Oxy$, cho $\vec{a} = (2; -3)$ và $\vec{b} = (-1; 4)$. Tọa độ của vectơ $\vec{u} = 2\vec{a} + \vec{b}$ là
\four{$(3; -2)$.}{$(5; 2)$.}{$(3; 2)$.}{$(1; 1)$.}

\item Cho tam giác $ABC$ có trọng tâm $G$. Với mọi điểm $M$ bất kì trong mặt phẳng, đẳng thức nào sau đây đúng?
\two{$\vec{MA} + \vec{MB} + \vec{MC} = 3\vec{MG}$.}{$\vec{MA} + \vec{MB} + \vec{MC} = 2\vec{MG}$.}{$\vec{MA} + \vec{MB} + \vec{MC} = 3\vec{GM}$.}{$\vec{MA} + \vec{MB} + \vec{MC} = \vec{MG}$.}

\item Trong mặt phẳng tọa độ $Oxy$, cho hai vectơ $\vec{x} = (m; 2)$ và $\vec{y} = (3; -6)$. Giá trị của $m$ để $\vec{x} \perp \vec{y}$ là
\four{$m = -4$.}{$m = 4$.}{$m = -1$.}{$m = 1$.}
\end{enumerate}

% ==================== PHẦN II: ĐÚNG / SAI ====================
\parttitle{PHẦN II. Câu trắc nghiệm đúng sai}
\partnote{Thí sinh trả lời từ câu 1 đến câu 2. Trong mỗi ý a), b), c), d) ở mỗi câu, thí sinh chọn đúng hoặc sai.}
\begin{enumerate}
\item Cho tam giác $ABC$ vuông tại $A$ có $AB = 6$, $AC = 8$ và $M$ là trung điểm cạnh $BC$ (xem hình vẽ bên).
\par\noindent\begin{minipage}[t]{0.74\linewidth}
\begin{enumerate}
    \item Độ dài cạnh huyền $BC = 10$.
    \item Độ dài vectơ tổng $|\vec{AB} + \vec{AC}| = 14$.
    \item Tích vô hướng $\vec{AB} \cdot \vec{AC} = 0$.
    \item Độ dài đường trung tuyến $AM = 5$.
\end{enumerate}
\end{minipage}\hfill
\begin{minipage}[t]{0.24\linewidth}\raggedleft\vspace{-8pt}
\begin{tikzpicture}[scale=0.58, >=Stealth]
  \coordinate (A) at (0,0);
  \coordinate (B) at (2.4,0);
  \coordinate (C) at (0,3.2);
  \coordinate (M) at (1.2,1.6);
  \fill[figFill] (A) -- (B) -- (C) -- cycle;
  \draw[thick, figEdge] (A) -- (B) -- (C) -- cycle;
  \draw[dashed, thick, figPrimary] (A) -- (M);
  \draw[figHighlight, semithick] (0,0.32) -- (0.32,0.32) -- (0.32,0);
  \foreach \pt in {A,B,C} \fill[figEdge] (\pt) circle (1.5pt);
  \fill[figHighlight] (M) circle (2pt);
  \node[below left, font=\scriptsize\sffamily, text=figEdge] at (A) {$A$};
  \node[below right, font=\scriptsize\sffamily, text=figEdge] at (B) {$B$};
  \node[above left, font=\scriptsize\sffamily, text=figEdge] at (C) {$C$};
  \node[above right, font=\scriptsize\sffamily, text=figHighlight] at (M) {$M$};
\end{tikzpicture}
\end{minipage}

\item Trong mặt phẳng tọa độ $Oxy$, cho ba điểm $A(1; 2)$, $B(4; 6)$ và $C(5; -1)$.
\begin{enumerate}
    \item Tọa độ vectơ $\vec{AB} = (3; 4)$.
    \item Độ dài đoạn thẳng $AB = 5$.
    \item Tọa độ trung điểm $I$ của $AC$ là $I(3; 1)$.
    \item Tích vô hướng $\vec{AB} \cdot \vec{AC} = 0$.
\end{enumerate}
\end{enumerate}

% ==================== PHẦN III: TRẢ LỜI NGẮN ====================
\parttitle{PHẦN III. Câu trắc nghiệm trả lời ngắn}
\partnote{Thí sinh trả lời từ câu 1 đến câu 2.}
\begin{enumerate}
\item Cho hình thoi $ABCD$ cạnh $10$ có góc $\widehat{BAD} = 60^\circ$. Tính độ dài vectơ $\vec{u} = \vec{AB} - \vec{AD}$. \answerbox

\item Trong mặt phẳng tọa độ $Oxy$, cho $\vec{a} = (3; 4)$ và $\vec{b} = (4; 3)$. Tính $\cos(\vec{a}, \vec{b})$ (viết dưới dạng số thập phân). \answerbox
\end{enumerate}

% ==================== PHẦN IV: TỰ LUẬN ====================
\parttitle{PHẦN IV. Câu tự luận}
\partnote{Thí sinh trình bày lời giải chi tiết từ Bài 1 đến Bài 2.}

\essayitem{Bài 1}{1,0 điểm} Cho tam giác $ABC$. Gọi $M$ là trung điểm $BC$ và $N$ là điểm thuộc cạnh $AC$ sao cho $NC = 2NA$.
\begin{enumerate}[label=\textsf{\color{slateBody}\alph*)}, leftmargin=0.6cm, itemsep=2pt]
    \item (0,5 điểm) Biểu diễn $\vec{AM}$ và $\vec{BN}$ theo hai vectơ $\vec{AB}$ và $\vec{AC}$.
    \item (0,5 điểm) Gọi $I$ là giao điểm của $AM$ và $BN$, đặt $\vec{AI} = k\vec{AM}$. Hãy tìm tỉ số $k$.
\end{enumerate}

\par\noindent\begin{minipage}[t]{0.72\linewidth}
\essayitem{Bài 2}{1,0 điểm} Một vật chịu tác dụng đồng thời của hai lực $\vec{F}_1, \vec{F}_2$ có độ lớn lần lượt là $30\text{ N}$ và $50\text{ N}$, hợp với nhau góc $60^\circ$ (như hình minh họa bên). Tính độ lớn hợp lực $\vec{F} = \vec{F}_1 + \vec{F}_2$.
\end{minipage}\hfill
\begin{minipage}[t]{0.26\linewidth}\raggedleft\vspace{-2pt}
\begin{tikzpicture}[scale=0.52, >=Stealth]
  \fill[figFill, rounded corners=1.5pt] (-0.35,-0.25) rectangle (0.35,0.25);
  \draw[thick, figEdge, rounded corners=1.5pt] (-0.35,-0.25) rectangle (0.35,0.25);
  \coordinate (O) at (0.35,0);
  \coordinate (F1) at (2.15,0);
  \coordinate (F2) at (1.85, 2.60);
  \draw[->, very thick, figPrimary] (O) -- (F1) node[below, font=\scriptsize, text=figPrimary] {$\vec{F}_1$};
  \draw[->, very thick, figSecondary] (O) -- (F2) node[above left, font=\scriptsize, text=figSecondary] {$\vec{F}_2$};
  \pic[draw=figHighlight, thick, angle radius=0.45cm, "$60^\circ$", angle eccentricity=1.68, font=\tiny, text=figHighlight] {angle = F1--O--F2};
\end{tikzpicture}
\end{minipage}

% ==================== TRANG ĐÁP ÁN & HƯỚNG DẪN GIẢI ====================
\newpage
\pagestyle{dapan}
\setcounter{page}{1}

{\color{navy}\hrule height 1.2pt}\vspace{6pt}
\begin{center}
  {\Large\sffamily\bfseries\color{ink} BẢNG ĐÁP ÁN TỔNG HỢP}\\[2pt]
  {\small\sffamily\color{muted} \examTitle\ — \examKind\ (Mã đề \examCode)}
\end{center}
\vspace{2pt}{\color{hairline}\hrule height 0.45pt}\vspace{6pt}

\noindent\textsf{\small\textbf{\color{ink}PHẦN I. Trắc nghiệm nhiều phương án lựa chọn}}
\vspace{3pt}

\begin{center}
\renewcommand{\arraystretch}{1.28}
\arrayrulecolor{hairline}
\begin{tabular}{|C{1.4cm}|C{1.4cm}|C{1.4cm}|C{1.4cm}|}
\hline
\rowcolor{surface}
\textsf{\textbf{\color{navy}1}} & \textsf{\textbf{\color{navy}2}} & \textsf{\textbf{\color{navy}3}} & \textsf{\textbf{\color{navy}4}} \\ \hline
\textbf{\color{ink}C} & \textbf{\color{ink}A} & \textbf{\color{ink}A} & \textbf{\color{ink}B} \\ \hline
\end{tabular}
\end{center}

\vspace{4pt}
\noindent\begin{minipage}[t]{0.52\linewidth}
\textsf{\small\textbf{\color{ink}PHẦN II. Trắc nghiệm đúng/sai}}\\[3pt]
\renewcommand{\arraystretch}{1.25}
\arrayrulecolor{hairline}
\begin{tabular}{|C{1.0cm}|C{0.75cm}|C{0.75cm}|C{0.75cm}|C{0.75cm}|}
\hline
\rowcolor{surface}
\textsf{\textbf{\color{navy}Câu}} & \textsf{\textbf{\color{navy}a)}} & \textsf{\textbf{\color{navy}b)}} & \textsf{\textbf{\color{navy}c)}} & \textsf{\textbf{\color{navy}d)}} \\ \hline
\textbf{\color{ink}1} & \textbf{\color{navy}Đ} & \color{muted}S & \textbf{\color{navy}Đ} & \textbf{\color{navy}Đ} \\ \hline
\textbf{\color{ink}2} & \textbf{\color{navy}Đ} & \textbf{\color{navy}Đ} & \color{muted}S & \textbf{\color{navy}Đ} \\ \hline
\end{tabular}
\end{minipage}\hfill
\begin{minipage}[t]{0.44\linewidth}
\textsf{\small\textbf{\color{ink}PHẦN III. Trả lời ngắn}}\\[3pt]
\renewcommand{\arraystretch}{1.25}
\arrayrulecolor{hairline}
\begin{tabular}{|C{1.0cm}|C{1.8cm}|}
\hline
\rowcolor{surface}
\textsf{\textbf{\color{navy}Câu}} & \textsf{\textbf{\color{navy}Đáp số}} \\ \hline
\textbf{\color{ink}1} & $10$ \\ \hline
\textbf{\color{ink}2} & $0{,}96$ \\ \hline
\end{tabular}
\end{minipage}

\vspace{8pt}
\parttitle{HƯỚNG DẪN GIẢI CHI TIẾT}
\vspace{2pt}

\begin{multicols}{2}
\raggedcolumns
\small
\arrayrulecolor{black}

\textsf{\textbf{\color{navy}PHẦN I. TRẮC NGHIỆM}}

\textbf{Câu 1.} Hình bình hành $ABCD$ có $\vec{AB} = \vec{DC}$. Chọn \textbf{C}.

\textbf{Câu 2.} $\vec{u} = 2(2; -3) + (-1; 4) = (3; -2)$. Chọn \textbf{A}.

\textbf{Câu 3.} Tính chất trọng tâm: $\vec{MA} + \vec{MB} + \vec{MC} = 3\vec{MG}$. Chọn \textbf{A}.

\textbf{Câu 4.} $\vec{x}\perp\vec{y} \Leftrightarrow 3m - 12 = 0 \Leftrightarrow m = 4$. Chọn \textbf{B}.

\textsf{\textbf{\color{navy}PHẦN II. ĐÚNG / SAI}}

\textbf{Câu 1.} a) $BC = 10$ (\textbf{Đ}). b) $|\vec{AB} + \vec{AC}| = 2AM = 10$ (\textbf{S}). c) $\vec{AB}\cdot\vec{AC} = 0$ (\textbf{Đ}). d) $AM = 5$ (\textbf{Đ}).

\textbf{Câu 2.} a) $\vec{AB} = (3; 4)$ (\textbf{Đ}). b) $AB = 5$ (\textbf{Đ}). c) $I(3; 0{,}5)$ (\textbf{S}). d) $\vec{AC} = (4; -3) \Rightarrow \vec{AB}\cdot\vec{AC} = 12 - 12 = 0$ (\textbf{Đ}).

\textsf{\textbf{\color{navy}PHẦN III. TRẢ LỜI NGẮN}}

\textbf{Câu 1.} $|\vec{AB} - \vec{AD}| = BD = \mathbf{10}$.

\textbf{Câu 2.} $\cos(\vec{a}, \vec{b}) = \dfrac{24}{25} = \mathbf{0{,}96}$.

\textsf{\textbf{\color{navy}PHẦN IV. TỰ LUẬN}}

\textbf{Bài 1.} a) $\vec{AM} = \dfrac{1}{2}\vec{AB} + \dfrac{1}{2}\vec{AC}$; $\vec{BN} = -\vec{AB} + \dfrac{1}{3}\vec{AC}$.\\
b) $\vec{BI} = \vec{AI} - \vec{AB} = \left(\dfrac{k}{2} - 1\right)\vec{AB} + \dfrac{k}{2}\vec{AC}$. Để $B, I, N$ thẳng hàng thì $\dfrac{k/2 - 1}{-1} = \dfrac{k/2}{1/3} \Leftrightarrow 1 - \dfrac{k}{2} = \dfrac{3k}{2} \Leftrightarrow k = \dfrac{1}{2}$.

\textbf{Bài 2.} $|\vec{F}| = \sqrt{30^2 + 50^2 + 2\cdot 30\cdot 50\cos 60^\circ} = 70\text{ N}$.

\end{multicols}
\end{document}
